\section{The first ten conditions at a coefficient row}

Fix an original member $D=w.D(Z)$ and one of its coefficient indices $i$.
Write $E/F$ for its CM extension, $c$ for complex conjugation,
$\ell=\ell_i$ and $\iota=\iota_i$ for the chosen complex embedding.
The member carries both a realization of its cuspidal labels and maps reading
their local parameters. B85U1 identifies the cuspidal labels with their
actual automorphic representations, injectively on classes of a fixed rank.
M84T1 identifies the period readings and M84T6 the character frame.
We use these three parts of the member in what follows.

% BEGIN S86A1
\paragraph{S86A1. Recovering the global list from its local factors.}
For a triple $(a,\pi,d)$ of positive ranks, expand its local
Speh parameter into the $d$ translates of $\pi$ by
\[
              (d-1)/2,\ (d-3)/2,\ \ldots,\ (1-d)/2.
\]
If two lists give equal parameters away from finitely many primes, the
standard unramified polynomials of the expanded lists agree there.
By Jacquet--Shalika uniqueness in A72F5 the two expanded multisets of
cuspidal classes with real norm twists agree. Each base cusp is unitary,
so its real twist exponent is recovered from the absolute value of the
central character (the norm exponent on the center is multiplied by
$a$). Removing the twist then determines its base cusp and hence its
original label. For each base cusp separately, let $m(s)$ count
occurrences at the nonnegative half-integer $s$. The number of strings
of length $2s+1$ is exactly $m(s)-m(s+1)$: all longer strings of
the same parity occur at both exponents, and the other parity occurs
at neither. This recovers ranks, labels and lengths with all repetitions.
For $d=1$ it gives ordinary cuspidal multiplicity one.
These are the two list comparisons in the selected row.

Split Frobenius elements away from any finite excluded set are dense in
$G_E$ when all lifts of places are allowed. Given a coset in a finite
Galois quotient over $E$, take a common normal closure over $F$. Lift
the requested element to the subgroup fixing $E$. Chebotarev over $F$
gives its conjugacy class outside the excluded primes, and choosing a
place above it gives this very lift. Its restriction to $E$ is trivial,
so the prime splits in $E/F$ and the residue norm at its lift in $E$
is unchanged. This meets the given coset. For the last clause of
$\mathrm{adm5}$ use the Brauer--Nesbitt theorem for semisimple
finite-dimensional representations with the same characteristic
polynomials on group elements, over the specified coefficient field.
In a basis its module isomorphism is the required simultaneous
conjugacy. This proves $\mathrm{an1},\mathrm{an2},\mathrm{a5}$.
% END S86A1

% BEGIN S86A2
\paragraph{S86A2. Characters, Satake roots and the local labels.}
For two rank-one labels with equal readings at almost all split places,
take the quotient $\xi$ of the unitary Hecke characters read by the member.
At split unramified places the inverse uniformizer values are precisely
the single roots of the given parameters. Hence $\xi$ is trivial there.
The finite-place Euler products for $L_E(s,\xi)$ and $\zeta_E(s)$ have a
quotient holomorphic and nonzero near $1$. Indeed at an inert prime
over $v$ their local logarithms have difference bounded by
$O((\mathrm Nv)^{-2\Re s})$ on compact sets in $\Re s>1/2$.
At the remaining finite exceptional primes the factors are regular and
nonzero near $1$ by unitarity. Consequently $L_E(s,\xi)$ has a pole
at $1$. The Hecke pole theorem (Tate's thesis, global zeta integral)
implies $\xi=1$, also for unitary characters with a norm component.
Injectivity of the character reading proves $\mathrm{HeckeSplitOne}$.

For each unramified positive-rank cuspidal label the Jacquet--Shalika
bound of A72F5 and its contragredient give
$q^{-1/2}<|\alpha|<q^{1/2}$ for all its Satake roots. The source
reads inverse roots, which satisfy the same pair of inequalities.
At all places a unitary character has bounded Weil image and zero
monodromy. This proves both clauses of $\mathrm{temp}$, even at
ramified places in rank one. The characters $\eta,\eta^{-1}$ are
conjugate self-dual: the defining restriction of $\eta$ to the real
subfield is its quadratic class character, trivial on norms, so
$\eta(x)\eta(c x)=1$. Their infinity multisets are the opposite
singletons $\{\pm1/2\}$ in the two embeddings, by the character
frame, and they and every $\mu_i$ are unitary. This gives
$\mathrm{TieAut}$.

At a place of a type prime $g$ the member reads the monic polynomial
whose roots lifted by its Teichm\"uller character are the eigenvalues of
the designated tame generator on the wild-invariant part of the
Weil space. A real norm twist is trivial on inertia, and this subspace
has dimension at most the full rank. Thus the label is unchanged by
the twist and its degree is bounded by the rank. This gives
$\mathrm{tieLab}$. Complete splitting of the type primes makes each
chosen place split over $F$, giving $\mathrm{gsp}$.
% END S86A2

% BEGIN S86A3
\paragraph{S86A3. Sums, twists and the original period readings.}
In the Weil--Deligne model the operations of direct sum, unramified
twist and tensor with a character are literal operations on
representations. Tensor distributes over direct sum, composes by
character multiplication and sends an inertia eigenvalue $z$ to
$\chi(\gamma)z$. The inertia polynomial is monic of degree the rank.
The trivial unramified line acts identically. Transporting through the
member's identification proves every identity in $\mathrm{TieLP}$.

At $w\mid\ell$, a de Rham representation becomes semistable over
a finite extension by the $p$-adic monodromy theorem (A72F2).
On its semistable period space $\ker N$ is the crystalline period
space there. Thus it is potentially crystalline exactly when it is
de Rham and $N=0$. It is crystalline over the original field exactly
when, in addition, the finite inertia descent action is trivial:
first descend the unramified constants by semilinear Frobenius
descent, while the ramified descent acts through inertia on the
period basis. Equivalently the resulting Weil--Deligne parameter is
unramified. Conversely a crystalline basis has these properties.
In a de Rham comparison each labelled filtered vector space has
dimension the rank, so the dimensions of its successive graded
quotients add to that rank, with repeated jumps counted separately.
The member's de Rham, jump and Weil--Deligne readings are those of
this comparison by its period frame and M84T1.

Away from $\ell$, write the inertia action in the form
$r(\gamma)\exp(t_\ell(\gamma)N)$ with $r$ of finite inertia
image. Frobenius semisimplification preserves this inertia and $N$.
An isomorphism of the parameter with a direct sum therefore gives
complementary inertia-stable spaces of precisely the two specified
dimensions. On a summand with $N=0$ a finite-order inertia matrix
in characteristic zero is semisimple. A characteristic polynomial
$(X-c)^m$ makes it scalar $c$ there, including after transport by
$\iota^{-1}$. Conjugation by a basis change transports each period
isomorphically. Direct sum takes direct sums of invariants and
filtrations. Tensor with a de Rham line shifts the labelled
filtration by that line's sole jump. Its inverse gives the converse
de Rham implication. These operations give the same sums and twists
of the associated Weil--Deligne parameters, away from $\ell$ by the
inertia formula and above $\ell$ by the semistable comparison after
one common finite extension. This supplies $\mathrm{TieGal}$ with its
de Rham guards. Together with S86A2 it proves $\mathrm{tie0}$.
% END S86A3

% BEGIN S86A4
\paragraph{S86A4. The cuspidal Galois attachment on this member.}
Fix a regular algebraic conjugate-self-dual cuspidal label $\pi$ of
rank $a\ge2$. Its realization from B85U1 has exactly these
properties, and its polarization character can be taken to be the
$a$th power of the quadratic character of $E/F$. This character is
trivial on norms and has the required sign $(-1)^a$ at infinity.
The contragredient identification gives the polarization of A72F5--8.
Take its single continuous semisimple $\ell$-adic attachment. It is
de Rham at \emph{every} $w\mid\ell$, and its jumps for an embedding
over $w$ are $(a-1)/2-z$, with $z$ running through the labelled
infinity multiset. We retain this unconditional de Rham conclusion
of the attachment theorem for use below.

The finite normalization is $\pi_w|\det|^{(1-a)/2}$. Both
Weil spaces have dimension $a$. At spherical places away from
$\ell$ the parameters agree. Under the source's additional inertia
test $\mathrm{Nc}$ the automorphic parameter is tame, has zero
monodromy and satisfies the stated repeated-root condition on the
chosen tame generators. An unramified twist preserves all three
properties. The member's cuspidal realization has good supercuspidal support
at every finite place (\(\mathrm{CGTrueT.good}\) in B85U1). Good
support is preserved by a norm twist. The good-support comparison
in A72F7 thus identifies the tested correspondence with the
classical one on these inputs over every completion. Apply the
polarized attachment and local comparison of A72F8. This gives
zero monodromy and equal inertia polynomials in the guarded case
away from $\ell$, and the zero-monodromy and inertia assertions of
$\mathrm{CuspA}$ in the guarded case above $\ell$. In the latter
case the domination of monodromy by a zero-monodromy parameter
forces zero monodromy on the Galois side as well. The equalities
are on the member's actual $\mathrm{wdOf}$ via its period
identification. When $a=1$, global reciprocity and the local character
comparison of A72F3 give the equality of the entire parameter at
every finite place and de Rhamness at every place above $\ell$.

Two further conclusions belong to this same field. For every regular
algebraic cuspidal label the conjugate infinity multiset is
$\{w_0-z\}$ by purity in A72F5. For every rank-one label $\chi$
and every real $s$ making all $z+s$ integral, the character
$\chi|\cdot|^s$ is algebraic. Its $\ell$-adic avatar in the source's
inverse-uniformizer convention has the indicated finite-part
parameter everywhere and labelled jumps $-(z+s)$ above $\ell$
by A72F1--3. This proves the entire $\mathrm{a2}$ field, including
these characters which were not assumed regular before the shift.
Finally the field called $\mathrm{g4p}$ in this version quantifies
the proposition $\mathrm{True}$ at each $i$ and is satisfied at once.
Together S86A1--4 establish all ten entries of the first part
$\mathrm{RowsAAlt}$ of $\mathrm{what}_2$ on the original member.
% END S86A4
